\section{成长经历：从县城到 ACM 班}

\subsection{童年与计算机的相遇}

陈天奇的中学在浙江省丽水市松阳县的松阳二中------他母亲是那里的高中老师，``当时觉得自己在自己地方上学也挺好的''。主持人问他童年对什么最感兴趣，他说最有趣的经历还是怎么开始接触计算机。

学习计算机的最初 motivation 是\textbf{想自己做电子游戏}------``学计算机的一个 motivation 就是能够自己创造游戏。当然这个 motivation 一直没有实现。''但正是这个对``创造''的渴望，引领他进入了编程世界，并一直延续到后来的科研中------``做机器学习也是这样的，创造智能。创造东西是一个原始的 motivation。''

\subsection{自学编程与信息竞赛}

初中时，他出于对做网页的向往找到了计算机启蒙老师何老师，得到一本 C 语言入门书。在没有竞赛教练、没有 GitHub 的年代，他通过网上的 Online Judge（OJ）和论坛（如``大龙树''）自学编程与信息竞赛。

\begin{knowledgebox}{自学者的``野路子''优势}
当时主流竞赛语言是 Pascal，但陈天奇因为入门时只学了 C，就一直用 C 打比赛。这种``野路子''反而培养了他不按套路出牌的习惯。高二时他出于兴趣，花了一个暑假从零写了一个 Pascal 到 C 的 transpiler（编译器雏形）。这段经历让他``不怕干事情''，奠定了日后从 first principle 出发解决问题的风格。
\end{knowledgebox}

高二参加 NOIP 初赛未过，旁听决赛拿了二等奖------当时看到结果``还蛮伤心的''。但第二年高三时拿到一等奖，而且高考完全没有落下，正常考入上海交通大学电院。主持人追问：``你当时在县城是不是一直全年级第一？''陈天奇笑答：``差不多算是吧。''

值得注意的是，学习编程的最初动机其实是想自己做电子游戏------这个 motivation 一直没有实现，但引领他进入了一个更广阔的世界。

\begin{warningbox}{关于``出生牛犊不怕虎''}
``出生牛犊不怕虎''是陈天奇在访谈中反复使用的词。从高中自学竞赛、到写 transpiler、到后来做 MXNet 挑战 TensorFlow、做 TVM 进入编译器领域，每一次他都是在``不知道会碰到什么问题''的情况下直接上手。这种心态的形成与他在小县城自学、没有权威告诉他``这个很难''有关------\textbf{没有见过天花板，反而不怕天花板}。
\end{warningbox}

\subsection{交大 ACM 班的历练}

2006 年，陈天奇通过高考考入上海交大电院，然后在宣讲会上认识了俞勇老师，精心准备材料后加入了 ACM 班。他回忆说，于（俞）老师后来跟他说``你是准备材料最仔细的一个学生''------这非常像后来准备 PhD 申请的 personal statement，``但当时并没有意识到这一点''。他作为一个``名不见经传的地方来的高考生''，能加入 ACM 班是``非常幸运的''。

ACM 班对他的影响体现在几个维度：

\begin{enumerate}
    \item \textbf{表达能力}：ACM 班的``教子讲坛''要求每个学生每学期选一个 topic 上台讲解，这培养了他在读博时直接 volunteer 做组会报告的自信。
    \item \textbf{``低调做人、高调做事''}：俞勇老师反复强调的价值观。``做一件事情要做到极致''成为他日后做 XGBoost 的核心驱动力。
    \item \textbf{编译器大作业}：ACM 班的编译原理课程设计要求从零搭建编译器，由第一届学长自发借鉴北美高校经验设计，一代一代传承。陈天奇认为这个训练``哪怕放到现在北美高校的标准来看都是一流的''。
    \item \textbf{体系结构大作业}：他在 ACM 班还挑战了在自己写的 CPU 上运行自己写的编译器。
\end{enumerate}

\begin{importantbox}{系统功底的积累}
ACM 班的编译器 + CPU 大作业经历，为陈天奇日后设计 TVM 编译器和 VTA 加速器打下了坚实基础。他回忆 TVM 开发时说：``因为有了原来那段经历，让我觉得不是很慌。''
\end{importantbox}

\begin{importantbox}{``低调做人、高调做事''的传承}
俞勇老师在 ACM 班反复强调的``低调做人、高调做事''深刻影响了陈天奇。他后来在职业生涯中接触各种成功人士后发现，真正能成事的人有两个特质：一是 ambitious，有锐气去挑战现状（challenge the status quo）；二是 humble，知道成功的因素``几乎不全部来自于你自己''，需要团队合作。这两者的结合------谦逊中保持锐气------成为他一贯的处事风格。
\end{importantbox}

\subsection{本章小结}

陈天奇的早期成长呈现出一个反复出现的模式：\textbf{在资源匮乏的环境中自学、不怕失败、从 first principle 出发解决问题}。从县城自学 OI，到高中写 transpiler，到 ACM 班写编译器跑在自己的 CPU 上，每一段经历都在为日后的大项目积累``不怕干事情''的勇气和系统设计的直觉。

